\subsection{预层的态射}

定义 3. --- 设 B 为一个拓扑空间，\(\mathcal{B}\) 为 B 的拓扑的一个基，\(\mathcal{F} = (\mathcal{F}(\mathrm{U}), f_{\mathrm{UV}})\) 与 \(\mathcal{G} = (\mathcal{G}(\mathrm{U}), g_{\mathrm{UV}})\) 为 B 上的、相对于 \(\mathcal{B}\) 的预层。从 \(\mathcal{F}\) 到 \(\mathcal{G}\) 中的预层态射是指一个从 \(\mathcal{F}\) 到 \(\mathcal{G}\) 的映射的投影系统。

换句话说 (E, III, p. 54)，从 \(\mathcal{F}\) 到 \(\mathcal{G}\) 中的预层态射是一个满足下列条件的族 \((\varphi_{\mathrm{U}})_{\mathrm{U}\in \mathcal{B}}\) ：

\((\mathrm{MPF}_1)\) 对每个属于 \(\mathcal{B}\) 的开集 U，\(\varphi_{\mathrm{U}}\) 是从 \(\mathcal{F}(\mathrm{U})\) 到 \(\mathcal{G}(\mathrm{U})\) 中的一个映射；

\((\mathrm{MPF}_2)\) 对每对开集 \((\mathrm{U},\mathrm{V})\) （属于 \(\mathcal{B}\) 且满足 \(\mathrm{U}\subset \mathrm{V}\) 者），有 \(\varphi_{\mathrm{U}}\circ f_{\mathrm{UV}} = g_{\mathrm{UV}}\circ \varphi_{\mathrm{V}}\) 。

当 \(\mathcal{F}\) 与 \(\mathcal{G}\) 是层时，从 \(\mathcal{F}\) 到 \(\mathcal{G}\) 中的预层态射也称为层态射。若 \(\mathcal{F}\) 与 \(\mathcal{G}\) 是 B 上的、相对于 \(\mathcal{B}\) 的预层，则从 \(\mathcal{F}\) 到 \(\mathcal{G}\) 中的预层态射构成一个集，记作 Mor( \(\mathcal{F};\mathcal{G}\) )。人们常不说“设 \(\varphi\) 为从 \(\mathcal{F}\) 到 \(\mathcal{G}\) 中的预层态射”，而说“设 \(\varphi\colon\mathcal{F}\to\mathcal{G}\) 为一个预层态射”。

设 \(\mathcal{F},\mathcal{G},\mathcal{H}\) 为 B 上的、相对于 \(\mathcal{B}\) 的预层，并设 \(\varphi \colon \mathcal{F}\to \mathcal{G}\) 与 \(\psi \colon \mathcal{G}\to \mathcal{H}\) 为预层态射。族 \((\psi_{\mathrm{U}}\circ \varphi_{\mathrm{U}})_{\mathrm{U}\in \mathcal{B}}\) 是从 \(\mathcal{F}\) 到 \(\mathcal{H}\) 中的预层态射，记作 \(\psi \circ \varphi\) 。族 \((\mathrm{Id}_{\mathcal{F}(\mathrm{U})})_{\mathrm{U}\in \mathcal{B}}\) 是从 \(\mathcal{F}\) 到其自身的预层态射，记作 \(\mathrm{Id}_{\mathcal{F}}\) 。

要使预层态射 \(\varphi = (\varphi_{\mathrm{U}})\colon \mathcal{F}\to \mathcal{G}\) 是一个同构，必须且只需对每个属于 \(\mathcal{B}\) 的开集 U，\(\varphi_{\mathrm{U}}\) 都是从 \(\mathcal{F}(\mathrm{U})\) 到 \(\mathcal{G}(\mathrm{U})\) 上的一个双射。这等价于说，存在一个预层态射 \(\psi \colon \mathcal{G}\to \mathcal{F}\) ，使得 \(\psi \circ \varphi = \operatorname{Id}_{\mathcal{F}}\) 且 \(\varphi \circ \psi = \operatorname{Id}_{\mathcal{G}}\) 。

设 \(\mathcal{F}\) 与 \(\mathcal{G}\) 为 B 上的、相对于 B 的拓扑的一个基 \(\mathcal{B}\) 的预层，设 B' 为 B 的一个开子集，并设 \(\mathcal{B}'\) 为 B' 的拓扑的一个基，使得 \(\mathcal{B}' \subset \mathcal{B}\) 。设 \(\varphi = (\varphi_{\mathrm{U}})_{\mathrm{U} \in \mathcal{B}}\) 为从 \(\mathcal{F}\) 到 \(\mathcal{G}\) 中的预层态射。那么 \((\varphi_{\mathrm{U}})_{\mathrm{U}\in \mathcal{B}^{\prime}}\) 是从 \(\mathcal{F}\mid \mathcal{B}'\) 到 \(\mathcal{G}\mid \mathcal{B}'\) 中的预层态射，记作 \(\varphi \mid \mathcal{B}'\) 。当 \(\mathcal{B}\) 是 B 的开子集全体且 \(\mathcal{B}'\) 是 \(\mathrm{B}'\) 的开子集全体时，\(\varphi \mid \mathcal{B}'\) 是从 \(\mathcal{F}\mid \mathrm{B}'\) 到 \(\mathcal{G}\mid \mathrm{B}'\) 中的预层态射，也记作 \(\varphi \mid \mathrm{B}'\) 。

例. --- 1) 设 B 为一个拓扑空间，(E, p) 与 (E', p') 为 B-空间，f: E → E' 为一个 B-态射。对 B 的每个开集 U，定义映射 \(f_{\mathrm{U}} \colon \mathcal{S}(\mathrm{U}; \mathrm{E}) \to \mathcal{S}(\mathrm{U}; \mathrm{E}')\) 为 \(f_{\mathrm{U}}(s) = f \circ s\) 。族 \(\underline{\mathcal{L}}(f) = (f_{\mathrm{U}})\) 是从 \(\underline{\mathcal{L}}(\mathrm{B}; \mathrm{E})\) 到 \(\underline{\mathcal{L}}(\mathrm{B}; \mathrm{E}')\) 中的预层态射。若 (E'', p'') 是一个 B-空间且 g: E' → E'' 是一个 B-态射，则有 \(\underline{\mathcal{L}}(g \circ f) = \underline{\mathcal{L}}(g) \circ \underline{\mathcal{L}}(f)\) 。

\begin{enumerate}
\def\labelenumi{\arabic{enumi})}
\setcounter{enumi}{1}
\item
  设 B 为一个拓扑空间，\(\mathcal{B}\) 为 B 的拓扑的一个基，\(\mathcal{F} = (\mathcal{F}(\mathrm{U}), f_{\mathrm{UV}})\) 为 B 上的、相对于 \(\mathcal{B}\) 的预层，\(\mathcal{L} = (\mathcal{L}(\mathrm{U}))\) 为 \(\mathcal{F}\) 的一个子预层。对每个开集 \(\mathrm{U} \in \mathcal{B}\) ，用 \(i_{\mathrm{U}}\) 表示从 \(\mathcal{L}(\mathrm{U})\) 到 \(\mathcal{F}(\mathrm{U})\) 中的典范单射。那么 \(i = (i_{\mathrm{U}})_{\mathrm{U} \in \mathcal{B}}\) 是从 \(\mathcal{L}\) 到 \(\mathcal{F}\) 中的预层态射。称 \(i\) 为从 \(\mathcal{L}\) 到 \(\mathcal{F}\) 中的典范态射。
\item
  设 B 为一个拓扑空间，\(\mathcal{B}\) 为 B 的拓扑的一个基，I 为一个集合。对每个 \(i\in\mathrm{I}\) ，设 \(\mathcal{F}_{i}=(\mathcal{F}_{i}(\mathrm{U}),f_{i,\mathrm{UV}})\) 为 B 上的、相对于 \(\mathcal{B}\) 的一个预层。用 \(\mathcal{F}\) 表示族 \((\mathcal{F}_{i})_{i\in\mathrm{I}}\) 的积预层。对每个开集 \(\mathrm{U}\in\mathcal{B}\) ，有 \(\mathcal{F}(\mathrm{U})=\prod_{i\in\mathrm{I}}\mathcal{F}_{i}(\mathrm{U})\) ；对每个 \(i\in\mathrm{I}\) ，用 \(\mathrm{pr}_{i,\mathrm{U}}\colon\mathcal{F}(\mathrm{U})\to\mathcal{F}_{i}(\mathrm{U})\) 表示指标 \(i\) 的投影。由预层 \(\mathcal{F}\) 的定义立得，族 \(\mathrm{pr}_{i}=(\mathrm{pr}_{i,\mathrm{U}})_{\mathrm{U}\in\mathcal{B}}\) 是从 \(\mathcal{F}\) 到 \(\mathcal{F}_{i}\) 中的预层态射。态射 \(\mathrm{pr}_{i}\) 称为指标 \(i\) 的投影态射。对每个预层 \(\mathcal{F}'\) （B 上的，相对于 \(\mathcal{B}\) 者）和每个族 \((\psi_{i})_{i\in\mathrm{I}}\) （其中 \(\psi_{i}\) 是从 \(\mathcal{F}'\) 到 \(\mathcal{F}_{i}\) 中的预层态射），存在唯一的预层态射 \(\psi\colon\mathcal{F}'\to\mathcal{F}\) ，使得对每个 \(i\in\mathrm{I}\) ，\(\mathrm{pr}_{i}\circ\psi=\psi_{i}\) 。
\item
  设 X 为 R 上的一个 \(C^{\infty}\) 类微分流形，并设 \(\underline{\mathcal{C}}^{\infty}(X;R)\) 为 X 上的 \(C^{\infty}\) 类数值函数的层。若 P 是 X 上的一个以 \(C^{\infty}\) 为系数的微分算子，则 P 到 X 的诸开集上的限制所成的族是层 \(\underline{\mathcal{C}}^{\infty}(X;R)\) 到其自身中的态射。可以证明，反之，层 \(\underline{\mathcal{C}}^{\infty}(X;R)\) 到其自身中的每个 R-线性态射都局部地具有这种形状 (I, p. 142, exerc. 3)。
\end{enumerate}

